\documentclass[10pt,a4paper]{amsart}
\usepackage{amsmath,amssymb,amsthm,graphicx}
\usepackage{verbatim}
\newenvironment{reference}{\comment}{\endcomment}
%\newenvironment{reference}{}{}
\newenvironment{slogan}{\comment}{\endcomment}
\newenvironment{history}{\comment}{\endcomment}

% For commutative diagrams we use Xy-pic
\usepackage[all]{xy}

% We use 2cell for 2-commutative diagrams.
\xyoption{2cell}
% Theorem environments.
%
\theoremstyle{plain}
\newtheorem{theorem}[subsection]{Theorem}
\newtheorem{proposition}[subsection]{Proposition}
\newtheorem{lemma}[subsection]{Lemma}

\theoremstyle{definition}
\newtheorem{definition}[subsection]{Definition}
\newtheorem{example}[subsection]{Example}
\newtheorem{exercise}[subsection]{Exercise}
\newtheorem{situation}[subsection]{Situation}

\theoremstyle{remark}
\newtheorem{remark}[subsection]{Remark}
\newtheorem{remarks}[subsection]{Remarks}

\numberwithin{equation}{subsection}

% Macros
%
\def\lim{\mathop{\mathrm{lim}}\nolimits}
\def\colim{\mathop{\mathrm{colim}}\nolimits}
\def\Spec{\mathop{\mathrm{Spec}}}
\def\Hom{\mathop{\mathrm{Hom}}\nolimits}
\def\Ext{\mathop{\mathrm{Ext}}\nolimits}
\def\SheafHom{\mathop{\mathcal{H}\!\mathit{om}}\nolimits}
\def\SheafExt{\mathop{\mathcal{E}\!\mathit{xt}}\nolimits}
\def\Sch{\mathit{Sch}}
\def\Mor{\mathop{\mathrm{Mor}}\nolimits}
\def\Ob{\mathop{\mathrm{Ob}}\nolimits}
\def\Sh{\mathop{\mathit{Sh}}\nolimits}
\def\NL{\mathop{N\!L}\nolimits}
\def\CH{\mathop{\mathrm{CH}}\nolimits}
\def\proetale{{pro\text{-}\acute{e}tale}}
\def\etale{{\acute{e}tale}}
\def\QCoh{\mathit{QCoh}}
\def\Ker{\mathop{\mathrm{Ker}}}
\def\Im{\mathop{\mathrm{Im}}}
\def\Coker{\mathop{\mathrm{Coker}}}
\def\Coim{\mathop{\mathrm{Coim}}}

% Boxtimes
%
\DeclareMathSymbol{\boxtimes}{\mathbin}{AMSa}{"02}

%
% Macros for moduli stacks/spaces
%
\def\QCohstack{\mathcal{QC}\!\mathit{oh}}
\def\Cohstack{\mathcal{C}\!\mathit{oh}}
\def\Spacesstack{\mathcal{S}\!\mathit{paces}}
\def\Quotfunctor{\mathrm{Quot}}
\def\Hilbfunctor{\mathrm{Hilb}}
\def\Curvesstack{\mathcal{C}\!\mathit{urves}}
\def\Polarizedstack{\mathcal{P}\!\mathit{olarized}}
\def\Complexesstack{\mathcal{C}\!\mathit{omplexes}}
% \Pic is the operator that assigns to X its picard group, usage \Pic(X)
% \Picardstack_{X/B} denotes the Picard stack of X over B
% \Picardfunctor_{X/B} denotes the Picard functor of X over B
\def\Pic{\mathop{\mathrm{Pic}}\nolimits}
\def\Picardstack{\mathcal{P}\!\mathit{ic}}
\def\Picardfunctor{\mathrm{Pic}}
\def\Deformationcategory{\mathcal{D}\!\mathit{ef}}
\usepackage{cleveref}
\crefname{originallemma}{original lemma}{original lemmas}
\Crefname{originallemma}{Original lemma}{Original lemmas}
\makeatletter
\newlabel{lemma-supported-on-support-pre}{{0GPF}{source}{}{sourcepointer-1}{}}
\newlabel{lemma-supported-on-support-pre@cref}{{[originallemma][0][]0GPF}{[1][0][]source}{}{}{}}
\makeatother
\title{The support subcategory equals the smallest Serre subcategory containing F closed under extensions and direct sums}
\author{The Stacks Project Authors}
\date{stacks-project:functors.tex:local-snapshot-20261012; GNU FDL 1.2 or later}
\begin{document}
\maketitle
\noindent\textit{Editorial note.} All delivered mathematical text is the author's own, copied verbatim from the pinned Stacks Project chapter file functors.tex; no statement or proof step was rewritten, reordered or removed.\par\medskip
\noindent\textit{Editorial note.} Span recomputation by label pairing (LANE\_BRIEF 10.2): in this chapter the author's proof environment follows the statement environment as a standalone block; the spec covers the complete environment pair (statement span, proof span) as recorded in the pool candidate.\par\medskip
\noindent\textit{Editorial note.} Occupation scan: combined claim\_key + statement\_anchor + generic-line scan over all registration files (active and withdrawn) found 0 hits for functors.tex (fu-occupied.json, 2026-10-14); consistent with the duty manager's six-check dispatch.\par\medskip
\ifdefined\hypertarget\hypertarget{sourcepointer-1}{}\fi
\noindent\textit{Source reference.} Original lemma 0GPF: sources/stacks-chow/tags.txt, source lines 20825--20825. Cross-chapter standard prerequisite cited by the author inside the counted span; it is used as a pointer and its own proof is not part of this item.\par\medskip
\begin{lemma}
\label{lemma-supported-on-support}
Let $X$ be a quasi-compact and quasi-separated scheme.
Let $\mathcal{F}$ be a finitely presented $\mathcal{O}_X$-module.
The following two subcategories of $\QCoh(\mathcal{O}_X)$ are equal
\begin{enumerate}
\item the full subcategory $\mathcal{A} \subset \QCoh(\mathcal{O}_X)$
whose objects are the quasi-coherent modules
whose support is (set theoretically) contained in $\text{Supp}(\mathcal{F})$,
\item the smallest Serre subcategory $\mathcal{B} \subset \QCoh(\mathcal{O}_X)$
containing
$\mathcal{F}$ closed under extensions and arbitrary direct sums.
\end{enumerate}
\end{lemma}

\begin{proof}
Observe that the statement makes sense as finitely presented
$\mathcal{O}_X$-modules are quasi-coherent.
Since $\mathcal{A}$ is a Serre subcategory closed under extensions and
direct sums and since $\mathcal{F}$ is an object of $\mathcal{A}$
we see that $\mathcal{B} \subset \mathcal{A}$. Thus it remains to
show that $\mathcal{A}$ is contained in $\mathcal{B}$.

\medskip\noindent
Let $\mathcal{E}$ be an object of $\mathcal{A}$. There exists
a maximal submodule $\mathcal{E}' \subset \mathcal{E}$
which is in $\mathcal{B}$. Namely, suppose
$\mathcal{E}_i \subset \mathcal{E}$, $i \in I$ is the set of
subobjects which are objects of $\mathcal{B}$. Then
$\bigoplus \mathcal{E}_i$ is in $\mathcal{B}$ and so is
$$
\mathcal{E}' = \Im(\bigoplus \mathcal{E}_i \longrightarrow \mathcal{E})
$$
This is clearly the maximal submodule we were looking for.

\medskip\noindent
Now suppose that we have a nonzero map
$\mathcal{G} \to \mathcal{E}/\mathcal{E}'$
with $\mathcal{G}$ in $\mathcal{B}$. Then
$\mathcal{G}' = \mathcal{E} \times_{\mathcal{E}/\mathcal{E}'} \mathcal{G}$
is in $\mathcal{B}$ as an extension of $\mathcal{E}'$ and $\mathcal{G}$.
Then the image $\mathcal{G}' \to \mathcal{E}$ would be strictly
bigger than $\mathcal{E}'$, contradicting the maximality of $\mathcal{E}'$.
Thus it suffices to show the claim in the following paragraph.

\medskip\noindent
Let $\mathcal{E}$ be an nonzero object of $\mathcal{A}$. We claim that there
is a nonzero map $\mathcal{G} \to \mathcal{E}$ with $\mathcal{G}$ in
$\mathcal{B}$. We will prove this by induction on the minimal
number $n$ of affine opens $U_i$ of $X$ such that
$\text{Supp}(\mathcal{E}) \subset U_1 \cup \ldots \cup U_n$.
Set $U = U_n$ and denote $j : U \to X$ the inclusion morphism.
Denote $\mathcal{E}' = \Im(\mathcal{E} \to j_*\mathcal{E}|_U)$.
Then the kernel $\mathcal{E}''$ of the surjection
$\mathcal{E} \to \mathcal{E}'$ has support contained in
$U_1 \cup \ldots \cup U_{n - 1}$. Thus if $\mathcal{E}''$ is nonzero,
then we win. In other words, we may assume that
$\mathcal{E} \subset j_*\mathcal{E}|_U$.
In particular, we see that $\mathcal{E}|_U$ is nonzero.
By Lemma \ref{lemma-supported-on-support-pre}
there exists a nonzero map $\mathcal{F}|_U \to \mathcal{E}|_U$.
This corresponds to a map
$$
\varphi : \mathcal{F} \longrightarrow j_*(\mathcal{E}|_U)
$$
whose restriction to $U$ is nonzero.
Setting $\mathcal{G} = \varphi^{-1}(\mathcal{E})$ we conclude.
\end{proof}

\end{document}
